% titel: Blatt 9, A6: Potenzreihe für ein Integral (Gauß-Integral)
% skript: Kap. 5.4 (Potenzreihen, gliedweise Integration; Beispiel 5.4.5)

% Teil I: Werkzeuge, kein Lösungsweg.

Eine Potenzreihe finden, die dasselbe leistet wie $\int_0^x \mathrm e^{-t^2/2}\,\dd t$
— also für alle $x$ konvergiert und diesen Integralwert darstellt — und den
Wert $\int_0^1\mathrm e^{-t^2/2}\,\dd t$ mit Python aproximieren.

\begin{werkzeugkasten}[Reihe der Exponentialfunktion]
Für alle $t\in\mathbb R$:
\begin{equation}
  \mathrm e^{u}=\sum_{k=0}^{\infty}\frac{u^k}{k!} .
\end{equation}
Mit $u=-t^2/2$ erhält man eine Reihe in geraden Potenzen von $t$, die für alle
$t$ (insbesondere $|x|\le 1$) konvergiert.

\textit{Skript:} Kap. 2.3, Beispiel 5.4.1.
\end{werkzeugkasten}

\begin{werkzeugkasten}[Gliedweise Integration einer Potenzreihe]
In ihrem Konvergenzintervall (hier $\mathbb R$) darf eine Potenzreihe
gliedweise integriert werden (Satz 5.4.2):
\begin{equation}
  \int_0^x \sum_{k=0}^{\infty}a_k t^k\,\dd t
  = \sum_{k=0}^{\infty}a_k\frac{x^{k+1}}{k+1}.
\end{equation}
Damit wird aus der Reihe für $\mathrm e^{-t^2/2}$ eine Potenzreihe
$\sum_n a_n x^n$ für den gewünschten Integralwert; da der Konvergenzradius
erhalten bleibt, konvergiert sie global.

\textit{Skript:} Satz 5.4.2, Beispiel 5.4.5, Punkt 3.
\end{werkzeugkasten}

\begin{werkzeugkasten}[Python-Approximation]
Die Reihe bis zu einem $N$ mit `sum` über eine Schleife aufsummieren, Glied
$k$ aus dem vorherigen Schritt rekursiv (Faktor pro $k$) — spart Laufzeit und
bleibt numerisch stabil. Einzelne Summanden
$(-1)^k \frac{x^{2k+1}}{(2k+1)\,2^k\,k!}$ schreiben.

\textit{Skript:} Anhang A.2 (Python-Befehle), S. 139.
\end{werkzeugkasten}

\begin{achtung}[Darauf achten]
\begin{itemize}
  \item Das Integral hat \emph{keine} elementare Stammfunktion; nur die
        Reihen-Darstellung liefert einen brauchbaren Weg.
  \item Beim gliedweisen Integrieren $t^{2k}\to\frac{x^{2k+1}}{2k+1}$ — alle
        $k$-abhängigen Faktoren sauber zusammensetzen.
  \item Die Antwort in a) muss eine \emph{Potenzreihe} $\sum_n a_n x^n$ sein,
        d. h. nach $x$ geordnete Koeffizienten $a_n$, nicht nur das Integral
        dahinschreiben.
  \item Bei $x=1$ alterniert die Reihe; abbrechen, wenn das nächste Glied klein
        genug ist.
\end{itemize}
\end{achtung}
